\label{ch:conclusion}

\section{Summary of Findings}
\label{sec:conclusion-summary}
\label{sec:research-contribution}

This paper is a complete specification of all the core sections that will be used in development of the project. It introduces a problem--solution view of how decentralized financing can be made adaptable to the institutional layers of our traditional financing system that is not flat, rather layered. Flat DeFi is hard to be made adaptable among the general mass, so our system provides a tier-based architecture for financial flow, with the added benefit of transparency with programmable smart contracts for the betterment of the client, and less control for the capitalistic organizations. The paper shows credible studies from the past that have contributed to the field before, and uses that evidence and those architectures as proof that a tier-based DeFi system can be built with better security layers that can contribute to the development finance sector globally. The core architecture stack follows mostly a traditional DeFi financial platform build, but the AI/ML layer integration and the tier-based layers are completely novel; no other paper or project shows all four-layer integration before. This project idea is enhanced by the existing systems that have two-tier-based banking systems.

A detailed plan has been introduced for how new clients can be onboarded and attracted to the platform, how clients will interact with the system, how banking authorities can approve loans, and how ML and AI layers can communicate in the middle with both clients and authorities to make this whole process of lending decentralized cryptocurrencies easier and more secure. The potential of such a decentralized system to be popular and adaptable is huge if implemented correctly and if the higher authorities show interest and invest heavily in this idea for the good of humanity. The purpose of this project is to push the idea of decentralized currencies forward for future generations to show more interest, and to support the building blocks for better financing systems in the future.

\section{Recommendations for Future Work}
\label{sec:future-work}

This section has been specified in the above chapters; it focuses on the potential of the project and the technical abilities and functional capabilities that can be improved and enhanced in the future. First priority is completing the MVT checklist of the development cycle for the defense of the project, and the rest are mostly added stretch to show how many features can be incorporated in the given time. A small summary of the future work:

Deployment on Polygon zkEVM Cardona, which remains experimental for production grade currently; as zkEVM may prove to be more stable, the platform will be deployed on Polygon zkEVM as it is more cost efficient for transactional activities.

Cross-tier features such as \texttt{InterBankLendingPool} and group lending, upward deposit facilities will be improved and made functional. For security and monitoring, multisig admin for World Bank operations and a live reserve transparency dashboard will be integrated. zkKYC (currently unstable globally), Certora proofs, federated learning, and GraphSAGE (quite complex to implement now) have been specified in the paper and will be more attainable with more time and effort; these have been moved to future work.

This field requires much research, as financial flow is heavily manipulated and misused globally by capitalists; transparency in capital flow is vital for equal access to resources and opportunity to live in harmony. The long-term goal of this project is to evolve this developed system to a more robust level, with more proper and efficient architecture and system design, a regulation-aware system while keeping transparency, equal access, and beneficial to humanity as core design principles.


